\ifdefined\gkver\else\def\gkver{student}\fi
\documentclass[\gkver]{gaokaozhenti}
\newcommand{\ccnum}[1]{{\fontspec{Noto Sans CJK JP}#1}}
\begin{document}

\chapter{2013年全国新课标Ⅰ}
%% paperType: 理综物理部分

\begin{enumerate}
\item
%% number: 14
%% paperName: 2013年全国新课标Ⅰ第 14 题 6 分
%% typeId: 1
%% type: 单选题
%% chapter: 直线运动
%% point: 匀变速直线运动
%% method: 无
%% score: 6
%% degree: 650
%% duplicateId: 0
%% body:
右图是伽利略 1604 年做斜面实验时的一页手稿照片，照片左上角的三列数据如下表。表中第二列是时间，第三列是物体沿斜面运动的距离，第一列是伽利略在分析实验数据时添加的。根据表中的数据，伽利略可以得出的结论是\xzanswer{C}

\onepicture[w=6cm, align=right]{figs/14.jpg}

\begin{center}
\begin{tabular}{|c|c|c|}
\hline
1 & 1 & 32 \\ \hline
4 & 2 & 130 \\ \hline
9 & 3 & 298 \\ \hline
16 & 4 & 526 \\ \hline
25 & 5 & 824 \\ \hline
36 & 6 & 1192 \\ \hline
49 & 7 & 1600 \\ \hline
64 & 8 & 2104 \\ \hline
\end{tabular}
\end{center}

\fourchoices[answer=C]
{物体具有惯性}
{斜面倾角一定时，加速度与质量无关}
{物体运动的距离与时间的平方成正比}
{物体运动的加速度与重力加速度成正比}

%% answer: C
%% memo:
\memoanswer{
伽利略对一个简单的加速度运动有两种猜测：一是物体的速度随位移均匀变化，另一个是物体的速度随时间均匀变化。他比较倾向后者，然后从数学上推理得出，如果物体的速度随时间均匀变化，则其位移将与时间的平方成正比（初速度为零时）。伽利略做这个实验的目的就是验证自己的后一个想法。正确选项 C。
}

\item
%% number: 15
%% paperName: 2013年全国新课标Ⅰ第 15 题 6 分
%% typeId: 1
%% type: 单选题
%% chapter: 电场
%% point: 电场叠加
%% method: 对称法
%% score: 6
%% degree: 650
%% duplicateId: 0
%% body:
如图，一半径为 $R$ 的圆盘上均匀分布着电荷量为 $Q$ 的电荷，在垂直于圆盘且过圆心 $c$ 的轴线上有 $a$、$b$、$d$ 三个点，$a$ 和 $b$、$b$ 和 $c$、$c$ 和 $d$ 间的距离均为 $R$，在 $a$ 点处有一电荷量为 $q$（$q>0$）的固定点电荷。已知 $b$ 点处的场强为零，则 $d$ 点处场强的大小为（$k$ 为静电力常量）\xzanswer{B}

\onepicture[w=4.5cm, align=right]{figs/15.png}

\fourchoices[answer=B]
{$k\frac{3q}{R^2}$}
{$k\frac{10q}{9R^2}$}
{$k\frac{Q+q}{R^2}$}
{$k\frac{9Q+q}{9R^2}$}

%% answer: B
%% memo:
\memoanswer{
由 $b$ 处的合场强为零可知圆盘在此处产生的场强与点电荷 $q$ 在此处产生的场强大小相等。$d$ 与 $b$ 关于圆盘对称，因此圆盘在 $d$ 处产生的场强与在 $b$ 处产生的场强大小相等。根据以上分析可知：$E_d=k\frac{q}{(3R)^2}+k\frac{q}{R^2}=k\frac{10q}{9R^2}$。正确选项 B。
}

\item
%% number: 16
%% paperName: 2013年全国新课标Ⅰ第 16 题 6 分
%% typeId: 1
%% type: 单选题
%% chapter: 电场
%% point: 带电粒子的直线运动
%% method: 无
%% score: 6
%% degree: 450
%% duplicateId: 0
%% body:
一水平放置的平行板电容器的两极板间距为 $d$，极板分别与电池两极相连，上极板中心有一小孔（小孔对电场的影响可忽略不计）。小孔正上方 $\frac{d}{2}$ 处的 P 点有一带电粒子，该粒子从静止开始下落，经过小孔进入电容器，并在下极板处（未与极板接触）返回。若将下极板向上平移 $\frac{d}{3}$，则从 P 点开始下落的相同粒子将\xzanswer{D}

\fourchoices[answer=D]
{打到下极板上}
{在下极板处返回}
{在距上极板 $\frac{d}{2}$ 处返回}
{在距上极板 $\frac{2d}{5}$ 处返回}

%% answer: D
%% memo:
\memoanswer{
第一个运动过程由动能定理可知：$mg\left(d+\frac{d}{2}\right)-qEd=0-0$；电容器保持与电源相连并且板间距减为原来的 $\frac{2}{3}$ 时，场强将由 $E$ 变为 $\frac{3}{2}E$，设粒子在距上极板 $x$ 的位置返回，则在此处时速度为零，由动能定理可知：$mg\left(x+\frac{d}{2}\right)-q\cdot\frac{3}{2}E\cdot x=0-0$。两式联立可得 $x=\frac{2}{5}d$。正确选项 D。
}

\item
%% number: 17
%% paperName: 2013年全国新课标Ⅰ第 17 题 6 分
%% typeId: 1
%% type: 单选题
%% chapter: 电磁感应
%% point: 切割磁感线电动势
%% method: 无
%% score: 6
%% degree: 450
%% duplicateId: 0
%% body:
如图，在水平面（纸面）内有三根相同的均匀金属棒 ab、ac 和 MN，其中 ab、ac 在 a 点接触，构成“V”字型导轨。空间存在垂直于纸面的均匀磁场。用力使 MN 向右匀速运动，从图示位置开始计时，运动中 MN 始终与 $\angle bac$ 的平分线垂直且和导轨保持良好接触。下列关于回路中电流 $i$ 与时间 $t$ 的关系图线，可能正确的是\xzanswer{A}

\onepicture[w=4.5cm, align=right]{figs/17a.png}

\fourchoices[answer=A, ispicture=true, h=2.3cm]
{figs/17b.png}
{figs/17c.png}
{figs/17d.png}
{figs/17e.png}

%% answer: A
%% memo:
\memoanswer{
在导体棒 MN 向右匀速运动过程中，由于其连入电路部分长度随时间线性增加，从而其电动势随时间线性增加。又由电阻定律可知，构成回路的三角形周长随时间线性增加，则其总电阻随时间线性增加。结合闭合电路欧姆定律可知，电流应恒定不变。正确选项 A。
}

\item
%% number: 18
%% paperName: 2013年全国新课标Ⅰ第 18 题 6 分
%% typeId: 1
%% type: 单选题
%% chapter: 磁场
%% point: 带电粒子在磁场中的运动
%% method: 无
%% score: 6
%% degree: 450
%% duplicateId: 0
%% body:
如图，半径为 $R$ 的圆是一圆柱形匀强磁场区域的横截面（纸面），磁感应强度大小为 $B$，方向垂直于纸面向外。一电荷量为 $q$（$q>0$）、质量为 $m$ 的粒子沿平行于直径 ab 的方向射入磁场区域，射入点与 ab 的距离为 $R/2$。已知粒子射出去的磁场与射入磁场时运动方向间的夹角为 $60^{\circ}$，则粒子的速率为（不计重力）\xzanswer{B}

\onepicture[w=3.6cm, align=right]{figs/18.png}

\fourchoices[answer=B]
{$\frac{qBR}{2m}$}
{$\frac{qBR}{m}$}
{$\frac{3qBR}{2m}$}
{$\frac{2qBR}{m}$}

%% answer: B
%% memo:
\memoanswer{
如图所示，粒子射出磁场与射入磁场时运动方向间的夹角为 $60^{\circ}$，则其运动轨迹所对的圆心角 $\angle NCM$ 也为 $60^{\circ}$。在三角形 $OHM$ 中，$\cos\angle HOM=\frac{1}{2}$，所以 $\angle HOM$ 为 $60^{\circ}$。由角边角定理可知，三角形 $OMN$ 与三角形 $CMN$ 全等，即圆周运动半径与磁场区域圆半径相等。由 $qvB=m\frac{v^2}{R}$ 可知 $v=\frac{qBR}{m}$。正确选项 B。

\onepicture[w=4cm]{figs/18_an.jpg}
}

\item
%% number: 19
%% paperName: 2013年全国新课标Ⅰ第 19 题 6 分
%% typeId: 2
%% type: 多选题
%% chapter: 直线运动
%% point: 运动图象
%% method: 无
%% score: 6
%% degree: 650
%% duplicateId: 0
%% body:
如图，直线 a 和曲线 b 分别是在平直公路上行驶的汽车 a 和 b 的位置-时间（$x-t$）图线。由图可知\xzanswer{BC}

\onepicture[w=4.5cm, align=right]{figs/19.png}

\fourchoices[answer=BC]
{在时刻 $t_1$，a 车追上 b 车}
{在时刻 $t_2$，a、b 两车运动方向相反}
{在 $t_1$ 到 $t_2$ 这段时间内，b 车的速率先减少后增加}
{在 $t_1$ 到 $t_2$ 这段时间内，b 车的速率一直比 a 车的大}

%% answer: BC
%% memo:
\memoanswer{
$t_1$ 时刻之前，b 在 a 的后面，而在此时间 b 追上 a，A 错。在 $x-t$ 图中，切线斜率表示物体的运动速度，其中斜率的绝对值表示速度的大小，斜率的正负表示运动方向，可知 BC 正确。正确选项 BC。
}

\item
%% number: 20
%% paperName: 2013年全国新课标Ⅰ第 20 题 6 分
%% typeId: 2
%% type: 多选题
%% chapter: 曲线运动
%% point: 人造地球卫星
%% method: 无
%% score: 6
%% degree: 650
%% duplicateId: 0
%% body:
2012 年 6 月 18 日，神州九号飞船与天宫一号目标飞行器在离地面 $343\Ukm$ 的近圆形轨道上成功进行了我国首次载人空间交会对接。对接轨道所处的空间存在极其稀薄的大气，下面说法正确的是\xzanswer{BC}

\fourchoices[answer=BC]
{为实现对接，两者运行速度的大小都应介于第一宇宙速度和第二宇宙速度之间}
{如不加干预，在运行一段时间后，天宫一号的动能可能会增加}
{如不加干预，天宫一号的轨道高度将缓慢降低}
{航天员在天宫一号中处于失重状态，说明航天员不受地球引力作用}

%% answer: BC
%% memo:
\memoanswer{
由于空气阻力的作用，其轨道半径越来越小，速度越来越大。失重并不是物体不受重力，而是对悬挂物的拉力或对支撑物的压力小于自身的重力，D 错误。在圆形轨道上绕地球做匀速圆周运动时，第一宇宙速度是其最大的环绕速度，A 错误。正确选项 BC。
}

\item
%% number: 21
%% paperName: 2013年全国新课标Ⅰ第 21 题 6 分
%% typeId: 2
%% type: 多选题
%% chapter: 运动和力
%% point: 牛顿定律的应用
%% method: 无
%% score: 6
%% degree: 650
%% duplicateId: 0
%% body:
2012 年 11 月，“歼 15”舰载机在“辽宁号”航空母舰上着舰成功。图（a）为利用阻拦系统让舰载机在飞行甲板上快速停止的原理示意图。飞机着舰并成功钩住阻拦索后，飞机的动力系统立即关闭，阻拦系统通过阻拦索对飞机施加一作用力，使飞机在甲板上短距离滑行后停止，某次降落，以飞机着舰为计时零点，飞机在 $t=0.4\Us$ 时恰好钩住阻拦索中间位置，其着舰到停止的速度-时间图线如图（b）所示。假如无阻拦索，飞机从着舰到停止需要的滑行距离约为 $1000\Um$。已知航母始终静止，重力加速度的大小为 $g$。则\xzanswer{AC}

\twopicture[numA={（a）}, numB={（b）}, wA=6cm, wB=7.5cm]
{figs/21a.png}
{figs/21b.png}

\fourchoices[answer=AC]
{从着舰到停止，飞机在甲板上滑行的距离约为无阻拦索时的 1/10}
{在 $0.4\Us\sim2.5\Us$ 时间内，阻拦索的张力几乎不随时间变化}
{在滑行过程中，飞行员所承受的加速度大小会超过 $2.5g$}
{在 $0.4\Us\sim2.5\Us$ 时间内，阻拦系统对飞机做功的功率几乎不变}

%% answer: AC
%% memo:
\memoanswer{
过 $(0,70)$ 和 $(3.0,0)$ 的直线与两坐标轴所围的三角形面积可替代曲线下的面积，其面积值为飞机在甲板上滑行的距离，A 正确。在 $0.4\Us\sim2.5\Us$ 时间内，飞机近似做匀变速运动，加速度不变，合力不变，阻拦索对飞机的合力不变，但由于阻拦索之间的夹角在变小，因此阻拦索中的力在变小，B 错误。由在 $0.4\Us\sim2.5\Us$ 时间内直线斜率可知飞机运动的最大加速度会超过 $2.5g$，C 正确。阻拦索对飞机的合力不变，但飞机的速度在减小，因此阻拦索系统对飞机做功的功率减小，D 错误。正确选项 AC。
}

\item
%% number: 22
%% paperName: 2013年全国新课标Ⅰ第 22 题 7 分
%% typeId: 4
%% type: 实验题
%% chapter: 直线运动
%% point: DIS测位移平均速度加速度
%% method: 无
%% score: 7
%% degree: 650
%% duplicateId: 0
%% body:
图（$a$）为测量物块与水平桌面之间动摩擦因数的实验装置示意图。实验步骤如下：

\begin{steps}
	\item
	用天平测量物块和遮光片的总质量 $M$、重物的质量 $m$；用游标卡尺测量遮光片的宽度 $d$；用米尺测量两光电门之间的距离 $s$；
	\item
	调整轻滑轮，使细线水平；
	\item
	让物块从光电门 A 的左侧由静止释放，用数字毫秒计分别测出遮光片经过光电门 A 和光电门 B 所用的时间 $\Delta t_A$ 和 $\Delta t_B$，求出加速度 $a$；
	\item
	多次重复上述测量，求 $a$ 的平均值 $\overline a$；
	\item
	根据上述实验数据求出动摩擦因数 $\mu$。
\end{steps}

回答下列问题：
\begin{enumerate}
	\item
	测量 $d$ 时，某次游标卡尺（主尺的最小分度为 $1\Umm$）的示数如图（b）所示，其读数为\tkanswer[1.4cm]{0.960}$\Ucm$。
	\item
	物块的加速度 $a$ 可用 $d$、$s$、$\Delta t_A$ 和 $\Delta t_B$ 表示为 $a=$\tkanswer[2.6cm]{$\frac{1}{2s}\left[\left(\frac{d}{\Delta t_A}\right)^2-\left(\frac{d}{\Delta t_B}\right)^2\right]$}。
	\item
	动摩擦因数 $\mu$ 可用 $M$、$m$、$\overline a$ 和重力加速度 $g$ 表示为 $\mu=$\tkanswer[2.6cm]{$\frac{mg-(M+m)\overline a}{Mg}$}。
	\item
	如果细线没有调整到水平，由此引起的误差属于\tkanswer[1.6cm]{系统误差}（填“偶然误差”或“系统误差”）。
\end{enumerate}

\twopicture[showlabel=false, wA=7.5cm, wB=4cm, align=right]
{figs/22a.png}
{figs/22b.png}

%% answer:
\jdanswer{
\begin{enumerate}
	\item
	$0.960\Ucm$；

	\item
	$\frac{1}{2s}\left[\left(\frac{d}{\Delta t_A}\right)^2-\left(\frac{d}{\Delta t_B}\right)^2\right]$；

	\item
	$\frac{mg-(M+m)\overline a}{Mg}$；

	\item
	系统误差。
\end{enumerate}
}
%% memo:
\memoanswer{
（1）游标卡尺的读数：$0.9\Ucm+12\times0.05\Umm=0.960\Ucm$。

（2）由于 $v_A=\frac{d}{\Delta t_A}$，$v_B=\frac{d}{\Delta t_B}$，由运动学公式知 $v_B^2-v_A^2=2as$，联立解得：$a=\frac{1}{2s}\left[\left(\frac{d}{\Delta t_A}\right)^2-\left(\frac{d}{\Delta t_B}\right)^2\right]$。

（3）由牛顿第二定律得：$mg-F_T=ma$；$F_T-\mu Mg=Ma$，联立解得 $\mu=\frac{mg-(M+m)a}{Mg}$。

（4）系统误差是由于仪器本身不精确，或实验方法粗略，或实验原理不完善而产生的。偶然误差是由各种偶然因素对实验者、测量仪器、被测物理量的影响而产生的。细线没有调整到水平引起的误差属于系统误差。
}

\item
%% number: 23
%% paperName: 2013年全国新课标Ⅰ第 23 题 8 分
%% typeId: 4
%% type: 实验题
%% chapter: 电路
%% point: 多用表的使用
%% method: 无
%% score: 8
%% degree: 650
%% duplicateId: 0
%% body:
某学生实验小组利用图（$a$）所示电路，测量多用电表内电池的电动势和电阻“$\times$1k”挡内部电路的总电阻。使用的器材有：

多用电表；

电压表：量程 $5\UV$，内阻十几千欧；

滑动变阻器：最大阻值 $5\UkO$；

导线若干。

回答下列问题：
\begin{enumerate}
	\item
	将多用电表挡位调到电阻“$\times$1k”挡，再将红表笔和黑表笔\tkanswer[1.2cm]{短接}，调零点。
	\item
	将图（$a$）中多用电表的红表笔和\tkanswer[1cm]{1}（填“1”或“2”）端相连，黑表笔连接另一端。
	\item
	将滑动变阻器的滑片调到适当位置，使多用电表的示数如图（$b$）所示，这时电压表的示数如图（$c$）所示。多用电表和电压表的读数分别为\tkanswer[1.4cm]{15.0}$\UkO$ 和\tkanswer[1.4cm]{3.60}$\UV$。
	\item
	调节滑动变阻器的滑片，使其接入电路的阻值为零。此时多用电表和电压表的读数分别为 $12.0\UkO$ 和 $4.00\UV$。从测量数据可知，电压表的内阻为\tkanswer[1.4cm]{12.0}$\UkO$。
	\item
	多用电表电阻挡内部电路可等效为由一个无内阻的电池、一个理想电流表和一个电阻串联而成的电路，如图（$d$）所示。根据前面的实验数据计算可得，此多用电表内电池的电动势为\tkanswer[1.4cm]{9.00}$\UV$，电阻“$\times$1k”挡内部电路的总电阻为\tkanswer[1.4cm]{15.0}$\UkO$。
\end{enumerate}

\onepicture[w=4.5cm, align=right]{figs/23a.png}

\onepicture[w=6cm, align=right]{figs/23b.png}

\onepicture[w=4.5cm, align=right]{figs/23c.png}

\onepicture[w=3.5cm, align=right]{figs/23d.png}

%% answer:
\jdanswer{
\begin{enumerate}
	\item
	短接；

	\item
	1；

	\item
	$15.0\UkO$，$3.60\UV$；

	\item
	$12.0\UkO$；

	\item
	$9.00\UV$，$15.0\UkO$。
\end{enumerate}
}
%% memo:
\memoanswer{
（1）欧姆表测电阻时，先选挡；然后进行欧姆调零，将红、黑表笔短接，调整欧姆调零旋钮，使指针指到欧姆表刻度的零位置；再测量、读数。

（2）对于多用电表欧姆挡，其电流从黑表笔流出，从红表笔流入，而电压表则要求电流从“$+$”接线柱流入，因此红表笔应和“1”端相连。

（3）多用电表欧姆挡的测量值等于表盘上读数乘以倍率，电阻“$\times$1k”读第一排，读数为 $15.0\times1\mathrm{k}=15.0\UkO$；直流电压 $5\UV$，最小分度为 $0.1\UV$，估读到分度值下一位，读数为 $3.60\UV$。

（4）调节滑动变阻器的滑片，使其接入电路的阻值为零，多用电表的读数即为电压表的内阻。

（5）由闭合电路欧姆定律得：$E=U+\frac{U}{R_V}(R_{\text{滑}}+r)$，有：$E=3.6+\frac{3.6}{12.0\times10^{3}}\left[(15-12)\times10^{3}+r\right]$，$E=4.0+\frac{4.0}{12.0\times10^{3}}r$。解得：$E=9.00\UV$，$r=15.0\UkO$。
}

\item
%% number: 24
%% paperName: 2013年全国新课标Ⅰ第 24 题 13 分
%% typeId: 6
%% type: 计算题
%% chapter: 直线运动
%% point: 相遇问题
%% method: 无
%% score: 13
%% degree: 650
%% duplicateId: 0
%% body:
水平桌面上有两个玩具车 A 和 B，两者用一轻质细橡皮筋相连，在橡皮筋上有一红色标记 R。在初始时橡皮筋处于拉直状态，A、B 和 R 分别位于直角坐标系中的 $(0,2l)$、$(0,-l)$ 和 $(0,0)$ 点。已知 A 从静止开始沿 $y$ 轴正向做加速度大小为 $a$ 的匀加速运动；B 平行于 $x$ 轴朝 $x$ 轴正向匀速运动。在两车此后运动的过程中，标记 R 在某时刻通过点 $(l,l)$。假定橡皮筋的伸长是均匀的，求 B 运动速度的大小。

%% answer:
\jdanswer{
$\frac{1}{4}\sqrt{6al}$
}
%% memo:
\memoanswer{
设 B 车的速度大小为 $v$。如图，标记 R 在时刻 $t$ 通过点 $K(l,l)$，此时 A、B 的位置分别为 H、G。由运动学公式，H 的纵坐标 $y_A$、G 的横坐标 $x_B$ 分别为 $y_A=2l+\frac{1}{2}at^2$\ccnum{①}，$x_B=vt$\ccnum{②}。

在开始运动时，R 到 A 和 B 的距离之比为 $2:1$，即 $OE:OF=2:1$。

由于橡皮筋的伸长是均匀的，在以后任一时刻 R 到 A 和 B 的距离之比都为 $2:1$。因此，在时刻 $t$ 有 $HK:KG=2:1$\ccnum{③}。

由于 $\triangle FGH$ 相似于 $\triangle IGK$，有 $HG:KG=x_B:(x_B-l)$\ccnum{④}，$HG:KG=(y_A+l):(2l)$\ccnum{⑤}。

由\ccnum{③}\ccnum{④}\ccnum{⑤}式得 $x_B=\frac{3}{2}l$\ccnum{⑥}，$y_A=5l$\ccnum{⑦}。

联立\ccnum{①}\ccnum{②}\ccnum{⑥}\ccnum{⑦}式得 $v=\frac{1}{4}\sqrt{6al}$\ccnum{⑧}。

\onepicture[w=5cm]{figs/24_an.jpg}
}

\item
%% number: 25
%% paperName: 2013年全国新课标Ⅰ第 25 题 19 分
%% typeId: 6
%% type: 计算题
%% chapter: 电磁感应
%% point: 电磁感应含容电路
%% method: 无
%% score: 19
%% degree: 450
%% duplicateId: 0
%% body:
如图，两条平行导轨所在平面与水平地面的夹角为 $\theta$，间距为 $L$。导轨上端接有一平行板电容器，电容为 $C$。导轨处于匀强磁场中，磁感应强度大小为 $B$，方向垂直于导轨平面。在导轨上放置一质量为 $m$ 的金属棒，棒可沿导轨下滑，且在下滑过程中保持与导轨垂直并良好接触。已知金属棒与导轨之间的动摩擦因数为 $\mu$，重力加速度大小为 $g$。忽略所有电阻。让金属棒从导轨上端由静止开始下滑，求：
\begin{enumerate}
	\item
	电容器极板上积累的电荷量与金属棒速度大小的关系；
	\item
	金属棒的速度大小随时间变化的关系。
\end{enumerate}

\onepicture[w=6cm, align=right]{figs/25.png}

%% answer:
\jdanswer{
\begin{enumerate}
	\item
	$Q=CBLv$；

	\item
	$v=\frac{m(\sin\theta-\mu\cos\theta)gt}{m+B^2L^2C}$。
\end{enumerate}
}
%% memo:
\memoanswer{
（1）设金属棒下滑的速度大小为 $v$，则感应电动势为 $E=BLv$\ccnum{①}。

平行板电容器两极板之间的电势差为 $U=E$\ccnum{②}。

设此时电容器极板上积累的电荷量为 $Q$，按定义有 $C=\frac{Q}{U}$\ccnum{③}。

联立\ccnum{①}\ccnum{②}\ccnum{③}式得 $Q=CBLv$\ccnum{④}。

（2）设金属棒的速度大小为 $v$ 时经历的时间为 $t$，通过金属棒的电流为 $i$，金属棒受到的磁场的作用力方向沿导轨向上，大小为 $f_1=BLi$\ccnum{⑤}。

设在时间间隔 $(t,t+\Delta t)$ 内流经金属棒的电荷量为 $\Delta Q$，按定义有 $i=\frac{\Delta Q}{\Delta t}$\ccnum{⑥}。

$\Delta Q$ 也是平行板电容器极板在时间间隔 $(t,t+\Delta t)$ 内增加的电荷量，由\ccnum{④}式得 $\Delta Q=CBL\Delta v$\ccnum{⑦}。

式中，$\Delta v$ 为金属棒的速度变化量，按定义有 $a=\frac{\Delta v}{\Delta t}$\ccnum{⑧}。

金属棒所受的摩擦力方向斜向上，大小为 $f_2=\mu N$\ccnum{⑨}。

式中，$N$ 是金属棒对于导轨的正压力的大小，有 $N=mg\cos\theta$\ccnum{⑩}。

金属棒在时刻 $t$ 的加速度方向沿斜面向下，设其大小为 $a$，根据牛顿第二定律有 $mg\sin\theta-f_1-f_2=ma$\ccnum{⑪}。

联立\ccnum{⑤}至\ccnum{⑪}式得 $a=\frac{m(\sin\theta-\mu\cos\theta)}{m+B^2L^2C}g$\ccnum{⑫}。

由\ccnum{⑫}式及题设可知，金属棒做初速度为零的匀加速运动。$t$ 时刻金属棒的速度大小为 $v=\frac{m(\sin\theta-\mu\cos\theta)}{m+B^2L^2C}gt$\ccnum{⑬}。
}

\item
%% number: 33
%% paperName: 2013年全国新课标Ⅰ第 33 题 9 分
%% typeId: 6
%% type: 计算题
%% chapter: 光学
%% point: 光的折射
%% method: 无
%% score: 9
%% degree: 650
%% duplicateId: 0
%% body:
\begin{enumerate}
	\item
	两个相距较远的分子仅在分子力作用下由静止开始运动，直至不再靠近。在此过程中，下列说法正确的是\xzanswer{BCE}（填正确答案标号。选对 1 个得 3 分，选对 2 个得 4 分，选对 3 个得 6 分。每选错 1 个扣 3 分，最低得分为 0 分）

	\fivechoices[answer=BCE]
	{分子力先增大，后一直减小}
	{分子力先做正功，后做负功}
	{分子动能先增大，后减小}
	{分子势能先增大，后减小}
	{分子势能和动能之和不变}
	\item
	如图，两个侧壁绝热、顶部和底部都导热的相同气缸直立放置，气缸底部和顶部均有细管连通，顶部的细管带有阀门 K。两气缸的容积均为 $V_0$，气缸中各有一个绝热活塞（质量不同，厚度可忽略）。开始时 K 关闭，两活塞下方和右活塞上方充有气体（可视为理想气体），压强分别为 $p_0$ 和 $p_0/3$；左活塞在气缸正中间，其上方为真空；右活塞上方气体体积为 $V_0/4$。现使气缸底与一恒温热源接触，平衡后左活塞升至气缸顶部，且与顶部刚好没有接触；然后打开 K，经过一段时间，重新达到平衡。已知外界温度为 $T_0$，不计活塞与气缸壁间的摩擦。求：
	\begin{enumerate}
		\item
		恒温热源的温度 $T$；
		\item
		重新达到平衡后左气缸中活塞上方气体的体积 $V_x$。
	\end{enumerate}

	\onepicture[w=5cm, align=right]{figs/33a.png}
\end{enumerate}

%% answer:
\jdanswer{
\begin{enumerate}
	\item
	BCE；

	\item
	$T=\frac{7}{5}T_0$，$V_x=\frac{1}{2}V_0$。
\end{enumerate}
}
%% memo:
\memoanswer{
（1）由分子力随分子间距离变化关系分析知，分子力先增大，然后减小，再增大，A 选项错误；分子从相距很远处开始运动，则 $r>r_0$ 时合力为引力，力和位移的夹角小于 $90^{\circ}$，分子力做正功，分子动能增大。$r<r_0$ 时合力为斥力，力和位移的夹角大于 $90^{\circ}$，分子力做负功，分子动能减小，BC 选项正确；由分子力做功与分子势能变化关系知，分子势能先减小，后增大，D 选项错误；分子仅在分子力作用下运动，只有分子力做功，分子势能和动能之间相互转化，分子势能和动能之和不变，E 选项正确。

\onepicture[w=5cm]{figs/33_an.png}

（2）与恒温热源接触后，在 K 未打开时，右活塞不动，两活塞下方的气体经历等压过程，由盖·吕萨克定律得 $\frac{T}{T_0}=\frac{7V_0/4}{5V_0/4}$\ccnum{①}，得 $T=\frac{7}{5}T_0$\ccnum{②}。

\onepicture[w=5cm]{figs/33_an2.jpg}

由初始状态的力学平衡条件可知，左活塞质量比右活塞的大。打开 K 后，左活塞下降至某一位置，右活塞必须升至气缸顶，才能满足力学平衡条件。

气缸顶部与外界接触，底部与恒温热源接触，两部分气体各自经历等温过程，设左活塞上方气体压强为 $p$，由玻意耳定律得 $pV_x=\frac{p}{3}\cdot\frac{V_0}{4}$\ccnum{③}，$(p+p_0)(2V_0-V_x)=p_0\cdot\frac{7}{4}V_0$\ccnum{④}。

联立\ccnum{③}\ccnum{④}式得 $6V_x^2-V_0V_x-V_0^2=0$\ccnum{⑤}，其解为 $V_x=\frac{1}{2}V_0$\ccnum{⑥}。另一解 $V_x=-\frac{1}{3}V_0$ 不合题意，舍去。

\onepicture[w=5cm]{figs/33_an3.jpg}
}

\item
%% number: 34
%% paperName: 2013年全国新课标Ⅰ第 34 题 6 分
%% typeId: 2
%% type: 多选题
%% chapter: 机械波
%% point: 波长频率波速
%% method: 无
%% score: 6
%% degree: 650
%% duplicateId: 0
%% body:
\begin{enumerate}
	\item
	如图，$a$、$b$、$c$、$d$ 是均匀媒质中 $x$ 轴上的四个质点，相邻两点的间距依次为 $2\Um$、$4\Um$ 和 $6\Um$。一列简谐横波以 $2\Ums$ 的波速沿 $x$ 轴正向传播，在 $t=0$ 时刻到达质点 $a$ 处，质点 $a$ 由平衡位置开始竖直向下运动，$t=3\Us$ 时 $a$ 第一次到达最高点。下列说法正确的是\xzanswer{ACD}（填正确答案标号。选对 1 个得 3 分，选对 2 个得 4 分，选对 3 个得 6 分。每选错 1 个扣 3 分，最低得分为 0 分）

	\onepicture[w=7cm]{figs/34a.png}

	\fivechoices[answer=ACD]
	{在 $t=6\Us$ 时刻波恰好传到质点 $d$ 处}
	{在 $t=5\Us$ 时刻质点 $c$ 恰好到达最高点}
	{质点 $b$ 开始振动后，其振动周期为 $4\Us$}
	{在 $4\Us<t<6\Us$ 的时间间隔内质点 $c$ 向上运动}
	{当质点 $d$ 向下运动时，质点 $b$ 一定向上运动}
	\item
	图示为一光导纤维（可简化为一长玻璃丝）的示意图，玻璃丝长为 $L$，折射率为 $n$，AB 代表端面。已知光在真空中的传播速度为 $c$。求：
	\begin{enumerate}
		\item
		为使光线能从玻璃丝的 AB 端面传播到另一端面，求光线在端面 AB 上的入射角应满足的条件；
		\item
		求光线从玻璃丝的 AB 端面传播到另一端面所需的最长时间。
	\end{enumerate}

	\onepicture[w=6cm, align=right]{figs/34b.png}
\end{enumerate}

%% answer: ACD
\jdanswer{
（2）$\sin i\leq\sqrt{n^2-1}$；$T_{\max}=\frac{Ln^2}{c}$。
}
%% memo:
\memoanswer{
（1）当 $t=6\Us$ 时，波传播的距离 $\Delta x=vt=2\times6\Um=12\Um$，A 选项正确；$t=0$ 时刻质点 $a$ 由平衡位置开始竖直向下运动，$t=3\Us$ 时 $a$ 第一次到达最高点，所以 $T=4\Us$，C 选项正确；质点 $c$ 恰好到达最高点的时间 $t=3\Us+6/2\Us=6\Us$，B 选项错误；$t=4\Us$ 时，质点 $c$ 位于波谷，D 项正确；由 $\lambda=vT$ 知 $\lambda=2\times4\Um=8\Um$，$bd=10\Um=\frac{5}{4}\lambda$，所以当质点 $d$ 从最高点向下运动时，质点 $b$ 从平衡位置向下运动，E 选项错误。

（2）（i）设光线在端面 AB 上 C 点（如图）的入射角为 $i$，折射角为 $r$，由折射定律有 $n=\frac{\sin i}{\sin r}$\ccnum{①}。

设该光线射向玻璃丝内壁 D 点的入射角为 $\alpha$，为了使该光线可在此光导纤维中传播，应有 $\alpha\geq\theta$\ccnum{②}，式中 $\theta$ 是光线在玻璃丝内发生全反射时的临界角，它满足 $\sin\theta=\frac{1}{n}$\ccnum{③}。

由几何关系得 $\alpha+r=90^{\circ}$\ccnum{④}。

由\ccnum{①}\ccnum{②}\ccnum{③}\ccnum{④}式得 $\sin i\leq\sqrt{n^2-1}$\ccnum{⑤}。

\onepicture[w=5cm]{figs/34_an.jpg}

（ii）光在玻璃丝中传播速度的大小为 $v=\frac{c}{n}$\ccnum{⑥}。

光速在玻璃丝中轴线方向的分量为 $v_x=v\sin\alpha$\ccnum{⑦}。

光线从玻璃丝端面 AB 传播到其另一端面所需时间为 $T=\frac{L}{v_x}$\ccnum{⑧}。

光线在玻璃丝中传播，在刚好发生全反射时，光线从端面 AB 传播到其另一面所需的时间最长，由\ccnum{②}\ccnum{③}\ccnum{⑥}\ccnum{⑦}\ccnum{⑧}式得 $T_{\max}=\frac{Ln^2}{c}$\ccnum{⑨}。
}

\item
%% number: 35
%% paperName: 2013年全国新课标Ⅰ第 35 题 9 分
%% typeId: 6
%% type: 计算题
%% chapter: 运动和力
%% point: 动量守恒定律
%% method: 无
%% score: 9
%% degree: 650
%% duplicateId: 0
%% body:
\begin{enumerate}
	\item
	一质子束入射到静止靶核 \ce{^{27}_{13}Al} 上，产生如下核反应：\ce{p + ^{27}_{13}Al -> X + n}，式中 p 代表质子，n 代表中子，X 代表核反应产生的新核。由反应式可知，新核 X 的质子数为\tkanswer[1.2cm]{14}，中子数为\tkanswer[1.2cm]{13}。
	\item
	在粗糙的水平桌面上有两个静止的木块 A 和 B，两者相距为 $d$。现给 A 一初速度，使 A 与 B 发生弹性正碰，碰撞时间极短。当两木块都停止运动后，相距仍然为 $d$。已知两木块与桌面之间的动摩擦因数均为 $\mu$，B 的质量为 A 的 2 倍，重力加速度大小为 $g$。求 A 的初速度的大小。
\end{enumerate}

%% answer:
\jdanswer{
\begin{enumerate}
	\item
	14，13；

	\item
	$v_0=\sqrt{\frac{28}{5}\mu gd}$。
\end{enumerate}
}
%% memo:
\memoanswer{
（1）根据核反应过程电荷数守恒和质量数守恒，新核 X 的质子数为 $1+13-0=14$，质量数为 $1+27-1=27$，所以中子数 $=27-14=13$。

（2）设在发生碰撞前的瞬间，木块 A 的速度大小为 $v$；在碰撞后的瞬间，A 和 B 的速度分别为 $v_1$ 和 $v_2$。在碰撞过程中，由能量和动量守恒定律，得
$\frac{1}{2}mv^2=\frac{1}{2}mv_1^2+\frac{1}{2}(2m)v_2^2$\ccnum{①}，
$mv=mv_1+2mv_2$\ccnum{②}，
式中，以碰撞前木块 A 的速度方向为正，由\ccnum{①}\ccnum{②}式得 $v_1=\frac{v_2}{2}$\ccnum{③}。

设碰撞后 A 和 B 的运动的距离分别为 $d_1$ 和 $d_2$，由动能定理得 $-\mu mgd_1=0-\frac{1}{2}mv_1^2$\ccnum{④}，$-\mu(2m)gd_2=0-\frac{1}{2}(2m)v_2^2$\ccnum{⑤}。

依题意有 $d=d_1+d_2$\ccnum{⑥}。

设 A 的初速度大小为 $v_0$，由动能定理得 $-\mu mgd=\frac{1}{2}mv^2-\frac{1}{2}mv_0^2$\ccnum{⑦}。

联立\ccnum{②}至\ccnum{⑦}式，得 $v_0=\sqrt{\frac{28}{5}\mu gd}$\ccnum{⑧}。
}

\end{enumerate}

\end{document}
